\subsection{Two-dimensional specializations of relation $(\mathrm{III}^\prime)$}

Write
\[
 B_4=\langle a,b,c\mid aba=bab,\ bcb=cbc,\ ac=ca\rangle.
\]
In the notation of \cite[\S10, Theorem~20]{H17}, the two inputs
to the harmonic auxiliary word $g_g=\beta_g^{-1}$ are
\[
 x_{45}=(ab)^3,\qquad x_{51}=(bc)^3.
\]
On the subgroup $D$, $g_g\in\widehat F(x,y)'$.

\begin{lemma}\label{lem:iii-prime-two-dimensional}
Let $k$ be any field and let $\rho:B_4\to\mathrm{GL}_2(k)$ be
a homomorphism.  Then $\rho((ab)^3)$ and $\rho((bc)^3)$ commute.
For a finite field $k$ and $g\in D$, it follows that
\[
 g_g\bigl(\rho((ab)^3),\rho((bc)^3)\bigr)=1.
\]
Thus such a direct matrix specialization of $(\mathrm{III}')$
does not retain its harmonic auxiliary-word term.
\end{lemma}
\begin{proof}
The elementary braid identity says that $(ab)^3$ is central in
$\langle a,b\rangle$, and $(bc)^3$ is central in
$\langle b,c\rangle$.  Both displayed images therefore belong
to the centralizer of $B=\rho(b)$.

If $B$ is not scalar, its centralizer in $M_2(k)$ is the
commutative algebra $k[B]$.  To see this directly, choose a vector
$v$ such that $v,Bv$ are a basis.  Such a vector exists for every
nonscalar two-by-two matrix.  An operator commuting with $B$ is
determined by its value on $v$, and the operators $\alpha I+\beta B$
give every possible such value.  In particular the two centralizer
elements commute.

If $B=tI$, the braid relation with $A=\rho(a)$ gives
$tA^2=t^2A$, and invertibility implies $A=tI$.  The other braid
relation similarly gives $\rho(c)=tI$.  All images then commute.

Over a finite field the representation extends continuously to
the profinite completion.  Evaluation of a commutator proword on
two commuting elements is one.  Applying this to $g_g$ proves the
second assertion.
\end{proof}

This lemma only excludes this direct way of retaining the
$g_g$ term.  It does not assert that $(\mathrm{III}')$ has no
consequences for the finite defects, nor that higher-dimensional
representations or comparisons of covering groups cannot provide
the required equality.  No vanishing or nonvanishing statement for
the two defects follows from it.
